% part: intuitionistic-logic
% chapter: introduction
% section: syntax

\documentclass[../../../include/open-logic-section]{subfiles}

\begin{document}

\olfileid{int}{int}{syn}

\olsection{نحو منطق شهودگرایانه}

نحو منطق شهودگرایانه همان نحو منطق گزاره‌ای است. در منطق گزاره‌ای
کلاسیک می‌توان برخی رابط‌ها را به‌کمک رابط‌های دیگر تعریف کرد؛ برای
نمونه، می‌توان~$!A \lif !B$ را با~$\lnot !A \lor !B$، یا~$!A \lor !B$
را با~$\lnot(\lnot !A \land \lnot !B)$ تعریف کرد. ازاین‌رو در
ارائه‌های منطق کلاسیک غالباً برخی رابط‌ها به‌عنوان اختصار این تعریف‌ها
معرفی می‌شوند. در منطق شهودگرایانه چنین نیست، جز در دو مورد:
$\lnot !A$ را می‌توان---و اغلب چنین می‌کنند---اختصار
$!A \lif \lfalse$ تعریف کرد. در آن صورت، البته، خود~$\lfalse$ نباید
تعریف شود!{} همچنین~$!A \liff !B$ را می‌توان، مانند منطق کلاسیک،
به‌صورت~$(!A \lif !B) \land (!B \lif !A)$ تعریف کرد.

!!^{formula}های منطق گزاره‌ای شهودگرایانه از
\emph{!!{propositional variable}ها} و ثابت گزاره‌ای~$\lfalse$ با
استفاده از \emph{رابط‌های منطقی} ساخته می‌شوند. داریم:

\begin{enumerate}
\item مجموعهٔ !!^a{denumerable}~$\PVar$ از !!{propositional variable}ها:
  $\Obj p_0$، $\Obj p_1$، \dots
\item ثابت گزاره‌ای برای !!{falsity}~$\lfalse$.
\item رابط‌های منطقی: $\land$ (عطف)، $\lor$ (فصل)، $\lif$
  (شرطی).
\item علائم سجاوندی: (، )، و ویرگول.
\end{enumerate}

\begin{defn}[فرمول]
\ollabel{defn:formulas} مجموعهٔ~$\Frm[L_0]$ از
\emph{!!{formula}های} منطق گزاره‌ای شهودگرایانه به‌طور استقرایی چنین
تعریف می‌شود:
\begin{enumerate}
\item $\lfalse$ یک !!{formula} اتمی است.

\item هر !!{propositional variable}~$\Obj p_i$ یک !!{formula} اتمی
است.

\item اگر~$!A$ و~$!B$، !!{formula} باشند، آنگاه~$(!A \land !B)$ یک
!!a{formula} است.

\item اگر~$!A$ و~$!B$، !!{formula} باشند، آنگاه~$(!A \lor !B)$ یک
!!a{formula} است.

\item اگر~$!A$ و~$!B$، !!{formula} باشند، آنگاه~$(!A \lif !B)$ یک
!!a{formula} است.

\tagitem{limitClause}{هیچ چیز دیگری !!a{formula} نیست.}{}
\end{enumerate}
\end{defn}

افزون بر رابط‌های اولیه‌ای که در بالا معرفی شدند، از نمادهای
\emph{تعریف‌شدهٔ} زیر نیز استفاده می‌کنیم: $\lnot$ (نفی) و~$\liff$
(!!{biconditional}). فرمول‌هایی که با عملگرهای تعریف‌شده ساخته می‌شوند
چنین فهمیده می‌شوند:
\begin{enumerate}
\item $\lnot !A$ اختصار~$!A \lif \lfalse$ است.

\item $!A \liff !B$ اختصار~$(!A \lif !B) \land (!B \lif !A)$ است.
\end{enumerate}

هرچند~$\lnot$ رسماً اختصار تلقی می‌شود، گاه در قواعد و بندهای تعریف‌ها
صریحاً با آن چنان رفتار می‌کنیم که گویی اولیه است. هدف اصلی این کار آن
است که بتوانیم مسئله‌های تمرینی را بیان کنیم.

\end{document}
